\section{Ordinary punctured quadratic action on a supplied normal-three germ}
\label{sec:r13-punctured-quadratic-action}

\paragraph{Scope and dependency convention.}
This section concerns the supplied original-field smooth integral
threefold model germ of BP and TC.  Its centre is denoted by $y$,
its original value group is $|K^\times|$, and
$F=\widetilde{\mathcal H(y)}$ has transcendence degree three over
the original finite residue field $k$.  The germ is ordinary, and
all later neighborhoods are ordinary.  The scalar presentation is
in the fixed finite continuous closed-cell PL/PS class of OR and
EN.  In particular, its coefficient functions have smooth
extensions near the closed cells of a finite rational complex.
No assertion about an arbitrary unsupplied threefold germ, higher
coupled relations, variable-PS closure, or GK tests is made here.

The inputs are the source-bound retained theorems OR, QA/QP,
EN/ES, BP, TC and FA, and the present-round FC, CC, SP, AH,
CS, SD/LG, UA and UF theorems.  The ordinary-locality statement
OL is also used.  Every use below preserves the complete actual
source and every original compact complementary smooth CLD test.
The improvement is an ordinary open-cover composition for the
action.  It does not construct one primitive on the whole
puncture or solve its global Killing descent problem.

Let $f$ be the original finite scalar tuple and write
\[
 P=\sum_\nu c_\nu[h_\nu][k_\nu],\qquad
 \kappa_2(P)=\sum_\nu c_\nu
             D_\delta h_\nu\wedge D_\delta k_\nu=0.
 \tag{PA1}
\]
Put
\[
 Z_P=\tfrac12\sum_\nu c_\nu
   (d'_Ph_\nu\wedge d''_Pk_\nu+
                         d'_Pk_\nu\wedge d''_Ph_\nu).
 \tag{PA2}
\]
The exact quadratic transgression gives $D[Z_P]=0$.  First adjoin
finitely many normalized residue generators of $F$ to $f$, so
the entry tuple has the full required residue rank.  Apply the
finite EN--AN/ES construction once, and freeze its final completed
tuple as $\tau$.  ES1 supplies the corrected conewise primitives
on every cone of this final complete ordinary image.  In
particular these include new cones created by its central unit
completion: the final paragraph of the ES proof pulls an EN
primitive from a containing old face, even for a rank-losing
projection, and then subtracts the common ambient correction.
There is no repeated central completion in this argument.
Later refinements only pull back these fixed data.  Its old
complete source is
\[
                  P_\tau=S_\tau=\{x:d_\tau(x)=3\}.
 \tag{PA3}
\]
The symbol $P_\tau$ in this section denotes a source set, and
is different from the formal relation $P$ in PA1.
Apply FA to a compact old source star in $P_\tau$, obtaining
$h=(\tau,g)$.  Its asserted injectivity is on $P_\tau$ only.
There is no assertion that $t_h$ embeds all of $S_h$.
Neither the old tuple $\tau$ nor the inventory $P_\tau$ will
be enlarged during the following cover argument.

\subsection{A genuine ordinary neighborhood at a top-cone point}

The next observation supplies the precise ordinary input at top
interiors.  It avoids assuming that a selected source cone is
already the complete ordinary moment image.

\begin{lemma}[Ordinary containment in a prepared top plane]
\label{lem:r13-top-plane-ordinary-containment}
Let $p$ be a closed triple point of a prepared TC model and let
$q_p(t_1,t_2,t_3)$ be its actual family, with all $t_a>0$.
Suppose the fixed old tuple is monomial there, with order matrix
$A$ of rank three.  Then at every sufficiently short positive
point $q=q_p(t)$ there is an ordinary neighborhood $V_q$ such
that
\[
                         t_\tau(V_q)\subset\operatorname{im}A.
 \tag{PA4}
\]
All positive radius vectors are included.  In particular no
rational-independence restriction is made.
\end{lemma}

\begin{proof}
Choose an integral basis $m_1,\ldots,m_b$ of the lattice of
rows annihilating $A$.  In the selected integral TC formal
chart each $\tau^{m_i}$ has normal exponent zero.  Its
constant residue $a_i$ belongs to the finite field
$\kappa(p)$ and is nonzero.  Let $H_i(T)\in K^\circ[T]$
lift the monic minimal polynomial of $a_i$ over $k$, with
unit constant term.  The formal expansion of
$H_i(\tau^{m_i})$ has constant coefficient divisible by
the original uniformizer.  Every coefficient not divisible
by that uniformizer has positive total normal degree.
At three radii strictly smaller than one, its Gauss norm
is therefore strictly smaller than one.  More explicitly,
after clearing a denominator with unit constant term, it
is bounded by the maximum of $|\pi|$ and the three radii.
The denominator has norm one.  The finite actual-to-formal
comparison of TC thus gives
\[
                  |H_i(\tau^{m_i})(q)|<1.
 \tag{PA5}
\]
Continuity makes all these inequalities hold on one ordinary
neighborhood of $q$.  For an element of norm less than one
the unit constant term of $H_i$ dominates, while for an
element of norm greater than one its monic leading term
dominates.  Hence PA5 forces $|\tau^{m_i}|=1$ there.
The common real kernel of the $m_i$ is $\operatorname{im}A$,
which proves PA4.  The argument uses a strict actual
inequality at each point, rather than a density argument.
\end{proof}

\begin{corollary}[Top-interior local action]
\label{cor:r13-top-interior-local-action}
At every positive point in the relative interior of a prepared
old top cone there is an ordinary neighborhood on which the
relation in PA1 has zero action for every allowed background.
\end{corollary}

\begin{proof}
EN supplies a continuous finite-PS covector on the closed old
cone, with the full prescribed strain, including coefficient
walls.  At a point in its relative interior choose a small open
set in its three-plane.  Extend the covector to an ambient
neighborhood by a linear retraction to that plane.  PA4 and
continuity give an ordinary neighborhood whose whole old moment
image lies in this open set of the plane.  The pulled covector
therefore has exactly the old energy strain at every ordinary
point and on every later face, also when the later tangent map
loses rank.  Add back the actual ambient exact ES correction.
QA then gives the claimed action equality.  Alternatively the
retained source-side regularization gives a fixed positive
control on each smaller local truncation.  A choice of primitive
on any other old cone is unnecessary here.
\end{proof}

\subsection{The fixed old inventory and its ordinary cover}

\begin{theorem}[Punctured action on the supplied normal-three germ]
\label{thm:r13-punctured-quadratic-action}
Under the supplied model and finite scalar hypotheses above,
there is an ordinary representative $U$ of $y$ such that for
every original symmetric middle-positive background $S$,
separately $d'$- and $d''$-closed, of legal bidegree $(p,p)$,
\[
                   \mathcal A_S(P)|_{U\setminus\{y\}}=0,
                         \qquad p=0,1.
 \tag{PA6}
\]
The equality holds on all original compact complementary smooth
CLD tests and all actual refinements.  Different members of the
ordinary cover used to prove PA6 may have different finite unit
tuples, covectors and positive controls.  The neighborhood $U$
and its cover are chosen for the fixed relation before choosing
$S$ or a later test.  Its absolute carrier is not changed.
\end{theorem}

\begin{proof}
We record the finite preparation and the ordinary coverage
separately, since they have different roles.

\emph{1. Fix the old source once.}
Use a compact source truncation in the analytic interior.  The
canonical source has a finite polyhedral structure, its moment
map has finite fibers, and its germ at $y$ is exhausted by the
finitely many closed TC cones.  The faithful coordinates $h$
embed precisely this old source truncation.  They are retained
whenever an actual point on it is compared with a model point.
Because its central moment fiber is finite, an ordinary
shrinking isolates its central point $y$ from the other points
of that fiber.  In faithful coordinates the germ is a finite
star with this single centre.  TC's uniform all-angle ordinary
convergence places a common short star in this shrinking.

Every later unit in the finite central geometric preparation
lists is defined near $y$ and restricts to a rational PL
function on $P_\tau$ by the canonical unit-log property.
Consequently a finite list of these units gives a common finite
subdivision of this same old star.  Equivalently, after putting
the list in a common enlarged source complex, $P_\tau$ is the
old subcomplex described by its canonical old-rank condition,
with that common subdivision.  The old set is not replaced by
the enlarged full source.  Actual coincident labels are compared
using $h$, not merely their old geometric image.  A test chosen
later may use units defined only on a punctured local member;
it is handled there by its actual common refinement, and is not
inserted into this central geometric inventory.

\emph{2. Cover the positive two-dimensional old pieces.}
There are finitely many labelled two-faces in the chosen old
subdivision.  FC prepares each by a proper residue curve and
two positive normal parameters, over the original field.  Apply
SP to these finitely many original curve blocks, retaining all
old coordinates and all marked points.  Apply the last faithful
separator after this Picard completion, with that completed tuple
as the old tuple for the FC source alternative.  SD then applies
to the entire final incident source: main branches, subdivisions,
new zero-old-column marks, and ordinary-plane satellite pieces
all remain present.  Its raw-kernel descent and LG correction
give an actual ambient local strain lift on every resulting
open positive two-dimensional sector.

There are only finitely many exceptional angular rays here.
Indeed SD's finite normal Newton lemma is applied to a fixed
finite list of final functions, once for each of the finitely
many old curve blocks.  It removes finitely many rational
one-dimensional tie walls and makes a finite regular toroidal
subdivision.  The all-radius CC comparison and AH projection
calibration justify SD on every open sector, including
commensurable parameters and all full sources over it.  The
ordinary neighborhood in SD is a neighborhood of the actual
point on the old face, with the complete moment image used in
the strain equality.  It is not a current restricted to the
curve or to a selected skeleton.

\emph{3. Cover the remaining finite ray families.}
Take all old one-faces, the finite Newton tie rays of step 2,
and all rays introduced by the finite geometric subdivisions.
This is still a finite list in the fixed old source star.
The regular monomial refinement in FC realizes each as a
divisor on a smooth projective SNC model.  The normalized old
$\tau$-character residues on this divisor generate a field of
transcendence degree two: at a free point of the old ray its
old graded rank is three and its value rank is one.  The
monomial valuation comparison identifies these residues with
the restrictions on the realized divisor.  Thus the DT Gauss
family itself belongs to $S_\tau$ at free radii, and at all
radii by continuity and closedness.  Faithfulness is applied
only after this old-source membership has been established.
Equality of the finite monomial coordinate norms and the
retained faithful coordinates then identifies its actual radial
family with the old family.  This is also the boundary
identification in TC; it holds at all radii by continuity after
identification at free parameters.

Apply UA and UF to each of these finitely many divisors and
to the pulled old energy.  They give a common short outer collar
and, for every compact interval in it, an ordinary neighborhood
of the entire completed core fiber carrying the transported
surface covector.  In particular every positive point of every
listed old ray has an ordinary neighborhood with a full strain
lift and zero action for every original $S$.  The local tuple
for one ray need not be used in the proof at another ray.
Choosing an interval around each radial value covers every
sufficiently short positive point.  Finiteness of the divisor
list permits a common initial shortening; compact subintervals
need not have one common primitive or coefficient bound.

\emph{4. Cover the top interiors.}
The remaining positive old points lie in relative interiors of
prepared top cones.  BP/TC and the finite monomial subdivision
make the fixed old functions monomial there.  The preceding
top-plane lemma and corollary give an ordinary all-background
action neighborhood at each such point, at every positive
radius vector.  Coefficient walls inside a top cone are covered
by its full EN covector; when a geometric subdivision records
them as faces, steps 2--3 supply the same local conclusion.

\emph{5. Cover the complement of the old source.}
At an ordinary point $x\notin P_\tau$ one has
\[
                     d_f(x)\le d_\tau(x)\le2,
                  \qquad s_f(x)\le2.
 \tag{PA7}
\]
OR therefore supplies an ordinary action-zero neighborhood of
$x$, with the original scalar tuple $f$ used for its rank
condition.  OR already retains arbitrary later tests and
refinements; their enlarged ranks need not satisfy PA7.

No recursively growing source inventory is required in this
step.  A new point revealed by units used in one local proof
is either outside $P_\tau$ and covered by PA7, or is a point
of the already classified old star.  In the latter case it
is a top, two-face, or ray point treated above.  If it is a
point of $P_\tau$ over the faithful old coordinates of a
specified core point, faithfulness on $P_\tau$ identifies it
with that actual point.  UF then includes it in the same
completed fiber.  Distinct new rays with other final coordinates
are not claimed to lie in UF's one product tube.

\emph{6. Use ordinary locality.}
After the finite common shortenings, every point of
$U\setminus\{y\}$ belongs to one of the ordinary neighborhoods
in steps 2--5.  This is a cover of actual points of the ordinary
puncture.  It is stronger than a cover of free parameters or of
an angular link.  For an original compact smooth test supported
in the puncture choose a finite subordinate original smooth
partition of unity.  Each localized test is annihilated by the
corresponding local action theorem.  OL gives PA6, using the
same fixed $S$ on each member.  Linearity of the action permits
their sum, and every local equality already allowed the
test's own actual common refinement.  This proves PA6 without
a global covector, a global fixed control, or a pushforward of
the arbitrary background.
\end{proof}

\paragraph{What this composition does not conclude.}
The finite cover argument is for the action on a punctured
ordinary domain.  It does not identify a full global primitive
for $Z_P$, remove the central point, prove a KC12 seminorm
certificate, or establish the general actual analytic relation
descent theorem.  In particular PK's real global Killing
obstruction does not contradict PA6, and PA6 does not remove
that obstruction from a demanded global primitive.

\subsection{The exact point-supported remainder}

\begin{lemma}[Separately closed order-zero currents at one point]
\label{lem:r13-point-supported-order-zero}
Let $X$ be an original allowed pure $n$-dimensional space and
let $T$ be a coefficient-Radon current of bidegree $(a,b)$,
supported at one actual point $y$.  If $d'T=0$ and $a<n$,
then $T=0$.  If $d''T=0$ and $b<n$, the same conclusion
holds.  Thus a separately closed such current is zero unless
$(a,b)=(n,n)$; in top degree it is $c\delta_y$ for a scalar
$c$.
\end{lemma}

\begin{proof}
For each fixed local smooth form $\xi$, the scalar functional
$u\mapsto T(u\xi)$ is a Radon measure, by the coefficient-
Radon property.  Its support is contained in $\{y\}$.
Consequently it depends only on $u(y)$.  In one common local
moment presentation, an arbitrary test can therefore be
replaced, for its pairing with $T$, by its form with constant
coefficients at $t(y)$, multiplied by a cutoff equal to one
near $y$.

Write the complementary first degree as $r=n-a>0$, and let
$\eta_0$ be this constant-coefficient form.  Put
$E=\sum_j(t_j-t_j(y))\partial/\partial t_j$, acting by
contraction on the first column only.  Then
\[
          \beta=\frac1r\iota'_E\eta_0,
                  \qquad d'\beta=\eta_0,
                  \qquad \beta(t(y))=0.
 \tag{PA8}
\]
For a compact smooth cutoff $\chi=1$ near $y$, closedness
gives $T(d'(\chi\beta))=0$.  The term involving $d'\chi$
vanishes on a neighborhood of $y$; equivalently every one of
its coefficient measures is supported where it is zero.
It follows that $T(\chi\eta_0)=0$, and hence $T$ annihilates
the original test.  This computation can be made in a common
refinement of any finite local test presentation, so it uses
all original tests.  Exchanging columns proves the other
assertion.  A top-degree coefficient-Radon current is a
signed Radon measure, and a measure supported at one point
is a scalar Dirac measure.
\end{proof}

\begin{proposition}[The supplied quadratic action has only a top Dirac remainder]
\label{prop:r13-quadratic-point-residual}
In the setting of Theorem~\ref{thm:r13-punctured-quadratic-action},
the action $\mathcal A_S(P)$ is a separately closed
coefficient-Radon current.  For $p=0$ it vanishes on $U$.
For $p=1$ there is a real scalar $c_S(P)$ with
\[
                       \mathcal A_S(P)=c_S(P)\delta_y
                              \quad\hbox{on }U.
 \tag{PA9}
\]
If $\chi$ is any original compact smooth scalar function equal
to one near $y$ and supported in $U$, then
$c_S(P)=\langle\mathcal A_S(P),\chi\rangle$.
The value is independent of $\chi$ and of a smaller ordinary
representative containing $y$.
\end{proposition}

\begin{proof}
Every scalar in PA1 is locally a finite continuous bounded
classical psh difference, by the retained R6 continuous-PS
DSH lemma (or directly by the finite convex-PL decomposition
in the PL case).  Expand the fixed finite signed expression.
The frozen bounded BT theorem for the fixed arbitrary $S$
makes every positive mixed term coefficient-Radon and separately
closed.  Its signed finite sum has the same two properties.
This is also stated explicitly in QA after QA11 for the
finite Hessian-product expressions.  It is not inferred for
an arbitrary distribution $D M_S(\alpha)$.

PA6 gives support in $\{y\}$.  The bidegree is $(p+2,p+2)$
on a threefold, so the preceding lemma proves the two cases.
The formula for $c_S(P)$ and its independence follow from
point support.  No regularization or local potential for the
arbitrary positive closed background has been assumed.
\end{proof}

\paragraph{A retained potential-background subclass.}
If, on an ordinary neighborhood of $y$, the degree-one
background is $S=D[v]$ for an actual locally bounded classical
psh potential, the retained BT symmetry already gives
\[
 \mathcal A_{D[v]}(P)
       =D\!\left(v\,\mathcal A_{[1]}(P)\right)=0.
 \tag{PA10}
\]
Indeed the actual delta zero in PA1 has zero ordinary current
for background $[1]$, and all factors in PA10 are within the
frozen signed bounded-potential calculus.  Finite signed sums
of such locally represented closed backgrounds satisfy the
same conclusion by linearity, when the final $S$ is allowed.
This is a direct retained corollary, not a construction of a
potential for an arbitrary $S$.  General middle-positive
separately closed degree-one backgrounds are not assumed to
have such a local representation.

\subsection{The remaining arithmetic locus does carry mass}

For completeness, ordinary OR locality alone gives the
following exact weaker support diagnostic in every dimension.
Let $G_f$ be the union of its action-zero neighborhoods at
points with $s_f\le2$.  Then $G_f$ is ordinary open and
\[
             \operatorname{supp}\mathcal A_S(P)
                     \subset U\setminus G_f
                     \subset\{x:s_f(x)\ge3\}.
 \tag{PA11}
\]
In dimension three the last condition forces
$d_f=3$ and $r_f=0$.  Thus all original scalar moment values
lie in $\mathbf Q\log|K^\times|$.  The arithmetic level set
need not be open or closed, and no measure-zero conclusion
follows from the countability of these levels.

Here is an exact normalization check.  On the original
$X=(\mathbf G_{m,K}^3)^{\mathrm{an}}$, put
$t_i=-\log|T_i|$ and $h_i=\max(0,t_i)$.  Then
\[
                     D h_1\wedge D h_2\wedge D h_3
                                  =\delta_{\xi_0},
 \tag{PA12}
\]
where $\xi_0$ is the original unit Gauss point and the mass
is exactly one.  The retained original torus calibration
assigns weight one to each primitive coordinate slope jump.
Equivalently, approximate each hinge by smooth convex functions
whose second derivatives are probability measures converging
to $\delta_0$.  The calibrated product pairing is
\[
  \int_{\mathbf R^3}\eta(\sigma(t))
              \prod_{i=1}^3 h_{i,\epsilon}''(t_i)\,dt,
 \tag{PA13}
\]
for every original compact smooth scalar test $\eta$.
It converges to $\eta(\xi_0)$.  Frozen local-uniform BT
convergence identifies this with the left side of PA12.
The coordinate torus lattice is primitive, so no lattice,
factorial, or absolute-value rescaling enters this mixed
product.  At $\xi_0$ the tuple $(T_1,T_2,T_3)$ has
$d_f=3$, $r_f=0$, and $s_f=3$.

This example is a positive BT measure with an atom precisely
on the rational-value locus.  It is not the action of an
actual quadratic zero.  Its role is to exclude the invalid
step from a dense lower-residue-rank locus, or from countably
many remaining moment levels, to the vanishing of PA9.
The remaining theorem needed for the supplied model is the
all-background scalar identity $c_S(P)=0$.

\paragraph{Source index.}
The retained order-zero and DSH inputs are QA,
\texttt{qa\_degree\_one\_and\_quadratic\_transgression.tex},
the paragraph following QA11 and QA18, and the R6 full source,
\texttt{lem:r6-continuous-pws-dsh}, together with its frozen
fixed-background bounded BT theorem.  The exact primitive
torus calibration is retained in R6's annulus unit-slope and
toric product sections and in GT3--GT5.  All geometric and
locality dependencies named at the beginning are source-bound
in the Round 12/13 indices.  The independent review of the
present composition must specifically check the fixed-old-source
bookkeeping and the ordinary coverage; a real angular cover
alone would not prove PA6.
